\documentclass[12pt,a4paper]{article}

\usepackage{fontspec}
\usepackage{polyglossia}
\setmainlanguage{russian}
\setotherlanguage{english}
\setmainfont{Times New Roman}
\setmonofont{Consolas}
\usepackage{geometry}
\usepackage{enumitem}
\usepackage{listings}
\usepackage{xcolor}

\geometry{margin=2.5cm}
\setlist{nosep}
\lstdefinestyle{grammar}{
  basicstyle=\ttfamily\small,
  columns=fullflexible,
  keepspaces=true,
  frame=single,
  framerule=0.2pt,
  xleftmargin=0.4cm,
  xrightmargin=0.4cm,
  breaklines=true
}

\title{Грамматика}
\date{}

\begin{document}

\maketitle

\section*{Терминалы}

Пробельные символы, кроме перевода строки после определения, считаются
разделителями и не входят в грамматику. Ключевое слово \texttt{let} не
считается переменной.

\begin{lstlisting}[style=grammar]
LET       ::= "let"
LAMBDA    ::= "\"
EQUALS    ::= "="
DOT       ::= "."
LPAREN    ::= "("
RPAREN    ::= ")"
NEWLINE   ::= "\n"
EOF       ::= end of input

VARIABLE  ::= LETTER { LETTER | DIGIT | "_" }
LETTER    ::= "a".."z" | "A".."Z"
DIGIT     ::= "0".."9"
\end{lstlisting}

\section*{Факторизованная EBNF-грамматика}

Грамматика описывает набор именованных определений и одно основное
лямбда-выражение. Лямбда-абстракции с несколькими параметрами разрешены:
\texttt{\textbackslash x y z. expr} означает
\texttt{\textbackslash x.\textbackslash y.\textbackslash z. expr}.

\begin{lstlisting}[style=grammar]
program              ::= { definition } expression EOF

definition           ::= LET VARIABLE EQUALS expression NEWLINE

expression           ::= abstraction
                       | application

abstraction          ::= LAMBDA variable-list DOT expression

variable-list        ::= VARIABLE { VARIABLE }

application          ::= operand { operand } [ abstraction ]

operand              ::= VARIABLE
                       | LPAREN expression RPAREN
\end{lstlisting}

В записи используются EBNF-операторы \texttt{\{...\}} и \texttt{[...]};
отдельных правил вида \texttt{A ::= epsilon} в грамматике нет. Для
парсера это означает обычные циклы и проверку следующего токена.


\section*{Примеры вывода}

\subsection*{Пример 1}

Вход:

\begin{lstlisting}[style=grammar]
let S = \x y z.x z (y z)
let K = \x y.x
S K K
\end{lstlisting}

Вывод по грамматике:

\begin{lstlisting}[style=grammar]
program
=> { definition } expression EOF
=> definition definition expression EOF
=> let S = expression NEWLINE let K = expression NEWLINE expression EOF
=> let S = abstraction NEWLINE let K = abstraction NEWLINE application EOF
=> let S = \ variable-list . expression NEWLINE let K = \ variable-list . expression NEWLINE S K K EOF
=> let S = \ x y z . application NEWLINE let K = \ x y . application NEWLINE S K K EOF
=> let S = \ x y z . x z (y z) NEWLINE let K = \ x y . x NEWLINE S K K EOF
\end{lstlisting}

После подстановки определений и бета-редукции нормальная форма:

\begin{lstlisting}[style=grammar]
\x.x
\end{lstlisting}

\subsection*{Пример 2}

Вход:

\begin{lstlisting}[style=grammar]
x \y.y
\end{lstlisting}

Вывод по грамматике:

\begin{lstlisting}[style=grammar]
program
=> { definition } expression EOF
=> expression EOF
=> application EOF
=> operand [ abstraction ] EOF
=> x abstraction EOF
=> x \ variable-list . expression EOF
=> x \ y . application EOF
=> x \ y . operand EOF
=> x \ y . y EOF
\end{lstlisting}

Это выражение читается как:

\begin{lstlisting}[style=grammar]
x (\y.y)
\end{lstlisting}

\subsection*{Пример 3}

Вход:

\begin{lstlisting}[style=grammar]
(\x.x) y
\end{lstlisting}

Вывод по грамматике:

\begin{lstlisting}[style=grammar]
program
=> { definition } expression EOF
=> expression EOF
=> application EOF
=> operand { operand } EOF
=> ( expression ) operand EOF
=> ( abstraction ) y EOF
=> ( \ variable-list . expression ) y EOF
=> ( \ x . application ) y EOF
=> ( \ x . x ) y EOF
\end{lstlisting}

\end{document}
